% Central Paper C content (journal-neutral). Edit here; journal wrappers in ../journals/ only \input these files.
% Every number traces to ../../FINDINGS.md.
\section{Data and methods}\label{sec:methods}

\subsection{Citation data and journals}
We use the parsed OpenCitations snapshot of January 2024 released through the Illinois Data Bank
\citep{arjomandi2024dataset}, with 1,962,840,983 citation links between integer-identified publications. Journal
(venue) assignments come from the OpenCitations Meta CSV dump \citep{ocmeta2024dump,massari2024meta}, joined on the
publication identifier; venue records are merged when they share an ISSN. Of 76.8 million journal articles,
72,720,703 obtain their venue from Meta, 388,394 from the venue field of the parsed file, and 3,695,762 have none. A
venue counts as a journal if its own Meta record has type \emph{journal} and it has an ISSN; this rule keeps 92,779
venues, of which 85,437 appear in the citation network. A stricter rule that also requires at least half of the
venue's items to be journal articles is used as a sensitivity check. Processing used DuckDB \citep{raasveldt2019duckdb}.

\subsection{JIF-window counts and journal scores}
For a citing year $t$, let $c(A\!\to\!B,t)$ be the number of citations from items of journal $A$ published in $t$ to
items of journal $B$ published in $t-1$ or $t-2$, the JIF window. We count citing items of any type but, following
Clarivate's practice of counting citations from indexed sources, only citing items that have a journal; including
citing items without a journal is a sensitivity analysis. For each recipient journal $B$ and year $t$ we compute the
numerator $n_B(t)=\sum_A c(A\!\to\!B,t)$, the self-citation share $c(B\!\to\!B,t)/n_B(t)$, a JIF proxy dividing $n_B(t)$
by $B$'s journal articles from $t-1$ and $t-2$, and the share of $n_B(t)$ from its largest external donor.

\subsection{Labels: Clarivate suppressions 2015--2023}
We rebuilt the suppression lists from primary Clarivate material: the ``Suppressed Titles 2007--2017'' spreadsheet
(data years 2015--2017: journal, year and type), the Journal Citation Reports help pages for the 2019, 2022, 2023 and
2024 releases (archived copies), the 2020 release page and a copy of the 2021 release list. The resulting table has 144
rows. There are 98 self-citation rows, 60 of them with a published self-citation share. There are 37 stacking rows:
17 named recipient--donor pairs with Clarivate's two stacking figures, 9 rows for the donors of those pairs, and 11
stacking journals without named partners (2015--2017). Three rows describe one expression of concern (a recipient
with three donors), and six ``suppression warnings'' are excluded. Journals were matched to our venues by ISSN taken
from Crossref journal records, with a documented alias or ISSN choice where titles were ambiguous; 131 of 138 names
matched. Table~\ref{tab:accounting} traces each analysis's denominator back to these rows.

We checked the table against Clarivate's current consolidated spreadsheet of suppressed titles (data years
2007--2024), retrieved in October 2026 from the Journal Citation Reports help centre. Every journal listed in that
spreadsheet for 2015--2023 is in our table. Nine self-citation rows taken from the lists published at the time (eight
for data year 2019, one for 2018) are absent from it; we keep them and report results without them as a sensitivity
check. For an out-of-sample test we took the suppressions for data years 2024 (from the same spreadsheet) and 2025
(from the help-centre page of the 2026 release) only after freezing our predictions (Section~\ref{sec:m-robust}).

\begin{table}[t]
\centering
\caption{From label rows to analysis denominators.}
\label{tab:accounting}
\small
\begin{tabular}{@{}lr@{}}
\toprule
Step & $n$ \\
\midrule
\multicolumn{2}{@{}l}{\emph{Self-citation fidelity (Section~\ref{sec:res-fidelity})}} \\
\quad self-citation rows & 98 \\
\quad with a published self-citation share & 60 \\
\quad matched to a journal & 59 \\
\quad evaluated ($\ge$20 numerator citations and citable items): references present / absent & 50 / 7 \\
\multicolumn{2}{@{}l}{\emph{Self-citation ranks}} \\
\quad matched self-citation rows & 95 \\
\quad ranked ($\ge$20 numerator citations and citable items) & 85 \\
\multicolumn{2}{@{}l}{\emph{Stacking fidelity}} \\
\quad named pairs with a published donor share & 18 \\
\quad both journals matched: donor references present / absent & 9 / 4 \\
\multicolumn{2}{@{}l}{\emph{CIDRE detection (Section~\ref{sec:res-cidre})}} \\
\quad named recipient--donor rows 2015--2023, incl.\ expression of concern & 20 \\
\quad both journals matched & 15 \\
\multicolumn{2}{@{}l}{\emph{Enrichment (Section~\ref{sec:res-variant})}} \\
\quad stacking-listed journals in the base, design cases excluded & 28 \\
\quad self-citation-listed journals in the base & 85 \\
\multicolumn{2}{@{}l}{\emph{Out-of-sample test (Section~\ref{sec:res-oos})}} \\
\quad journals newly listed for stacking, data years 2024--2025 & 8 \\
\quad in the evaluated population & 6 \\
\bottomrule
\end{tabular}
\end{table}

\subsection{Fidelity measures}
Clarivate reports, for self-citation cases, the percentage of self-citations in the JIF numerator and, for stacking
pairs, the donor's percentage of the recipient's JIF numerator and the percentage exchange to JIF years, defined as
``the proportion of all citations from donor to recipient (all years) that reference the two years considered in the
Journal Impact Factor calculation''. We compute the same quantities; for the exchange share our denominator counts
cited years before $t$, so the estimate is an upper bound. We fixed in advance that a figure counts as reproduced if it lies within $\pm10$ percentage points,
and that a listed journal ``ranks'' if it is in the top 1\% of its year on the relevant share.

\subsection{CIDRE}
We implemented CIDRE as described by \citet{kojaku2021cidre}, wrapping the authors' package and closing
three gaps between package and paper: condition 1 (more than half of the $A\!\to\!B$ citations in year $t$ go to the
JIF window), the minimum of $\theta_w=50$ within-group citations, and a compatibility fix. An edge is excessive if it
meets condition 1 and its weight is significant under a Poisson test against the expected weight from a
degree-corrected stochastic block model, with Benjamini--Hochberg correction at $\alpha=0.01$
\citep{benjamini1995}. A journal's donor (recipient) score is the share of its outgoing (incoming) citations that are
excessive within the current group; journals whose donor and recipient scores are both below $\theta=0.15$ are
removed iteratively, and the remaining weakly connected components are the groups. The block model was fitted once on
the network aggregated over all years with graph-tool's nonparametric inference
\citep{peixoto2017nonparametric,peixoto2014graphtool}. Detection follows the paper's rule: a Clarivate group is detected
if at least half of a CIDRE group's journals and at least two journals are shared.

The block-model fit is stochastic. We therefore fitted seven partitions of the main journal network (five seeds, a
repeat of the configuration seed, and a single-threaded fit) and one of the stricter journal rule, and ran CIDRE on
every year from 2015 to 2023 with each.

\subsection{A recipient-centric reading of CIDRE}
Clarivate's stacking criterion is recipient-centric: a recipient that depends on one donor for a large share of its
numerator. CIDRE's excessive-edge test can be read the same way, without the group step. For journal $j$ in year
$t$, let $e_j$ be the number of excessive citations it receives (self-citations excluded) and $s_j$ its in-strength.
We flag $j$ if $e_j/s_j \ge \theta$ and $e_j \ge \theta_w$, with CIDRE's own $\theta=0.15$ and $\theta_w=50$. This
variant was designed after we diagnosed CIDRE's misses on two 2023 cases (Section~\ref{sec:res-cidre}); those two
journals are therefore excluded from every evaluation in which they could favour it. Because flags vary between
partitions, we also report consensus flags present in at least $k$ of the seven partitions. We evaluate the variant by
enrichment: the share of listed journals ever flagged in 2015--2023 compared with the share of all journals with at
least 20 numerator citations in some year, with a one-sided binomial test.

\subsection{Baselines, size matching and an out-of-sample test}\label{sec:m-robust}
These checks were added after an internal review of a first draft; their protocols were recorded in the lab notebook
before any of their results were computed.

\emph{Baselines.} A plain rule without a null model flags journal $j$ in year $t$ if its largest external donor
provides at least a share $X$ of its numerator (and the numerator is at least 20). We set $X$ on all journals, without
labels, so that the rule flags the same share of journals as the variant; as a second version we apply CIDRE's two
thresholds to the plain share (share $\ge0.15$ and at least 50 citations from the top donor). To isolate the role of
the block model, we also run the variant's test against a single block, that is, a configuration-model null.

\emph{CIDRE with an asymmetric threshold.} We rerun CIDRE with the recipient threshold kept at $0.15$ and the donor
threshold lowered to 0.10, 0.05, 0.02 and 0.01, so that a journal is removed only if its donor score is below the
donor threshold and its recipient score below $0.15$.

\emph{Size.} Stacking-listed journals tend to be small. We therefore compare them with journals of similar size,
measured by the largest yearly numerator: by a Mantel--Haenszel analysis over size deciles, and against ten
never-listed controls per listed journal of nearest size, drawn from the same block where possible.

\emph{Out of sample.} Before opening Clarivate's lists for data years 2024 and 2025, we froze the set of journals
flagged in 2015--2023 (a file with a recorded checksum). Among journals with at least 20 numerator citations in some
year and not listed in 2015--2023, we then test whether flagged journals were newly listed for stacking more often
than others (one-sided Fisher exact test). Our citation data end with data year 2023, so this tests whether a flag
anticipates a later listing.

\subsection{Paper-level anatomy and evidence check}
For flagged recipients and their largest excessive donor, and for named Clarivate pairs, we collect the citing papers
behind the pair's window citations and compute the citations per citing paper, the Gini coefficient of citations
over citing papers, and the shares coming from the single top citing paper and from the top 10\% of citing papers. We
compare titles of citing and cited papers by TF--IDF cosine similarity. Each pair receives up to three controls from
the same year with 0.8 to 1.25 times its window citations, both journals in the journal rule, and a recipient never
flagged. As a structured evidence check, fixed before the sample was drawn, we record three public indicators per
pair: either journal appears in Clarivate's lists; recipient and donor share a Crossref publisher; or the top 10\% of
citing papers carry at least half of the pair's citations.

\subsection{Reproducibility}
Every analysis is a script that writes a run folder with its parameters, input and output checksums and log, and the
decisions and pass criteria were recorded in a lab notebook before the corresponding runs. We distinguish
pre-specified from post hoc analyses throughout.

